\documentclass[11pt,a4paper]{article}

\usepackage[margin=1in]{geometry}
\usepackage{amsmath,amssymb,amsthm}
\usepackage{booktabs}
\usepackage{graphicx}
\usepackage{subcaption}
\usepackage{hyperref}
\usepackage{xcolor}
\usepackage{microtype}
\usepackage{parskip}
\usepackage{cleveref}
\usepackage{float}
\usepackage{multirow}
\usepackage{array}
\usepackage{caption}


\newcommand{\R}{\mathbb{R}}
\newcommand{\frob}[1]{\left\lVert #1 \right\rVert_F}
\newcommand{\muval}{\mu}
\newcommand{\lams}{\lambda_{\text{step}}}
\newcommand{\Perm}{\mathcal{P}}
\newcommand{\Birk}{\mathcal{B}}


\title{%
  \textbf{[WORK IN PROGRESS] \\
  An Investigation of Optimal Transport based Graph Alignment Algorithms\\}
}

\author{Archana Narayanan}
\date{March 2026}


\begin{document}
\maketitle


\begin{abstract}
Graph alignment has broad applications across scientific domains, yet current state-of-the-art methods remain far from ground-truth solutions.
We investigate this gap by comparing two leading graph matching approximation algorithms: FUGAL, an optimal-transport–based method, and the widely used FAQ.
We evaluate the impact of four algorithmic factors: initialization, learning-rate scheduling, permutation matrix projection, and additional terms in the objective function.
Results indicate that performance depends on both dataset characteristics and parameter settings, rather than a single optimal configuration. Warm-start initialization greatly improves FAQ but introduces impractical runtime costs, whereas FUGAL matches this performance without costly initialization in most cases.
Our study demonstrates that significant accuracy gains can be achieved through principled algorithmic tuning alone, highlighting unexplored opportunities to approach ground-truth alignment without added computational complexity.





\end{abstract}


\section{Introduction}
\label{sec:intro}

 Graph alignment seeks to find the best correspondence between the nodes of two graphs.

 FUGAL (Feature-Fortified Unrestricted Graph Alignment) is a state-of-the-art graph alignment algorithm that has been shown to surpass
  prior methods in accuracy across standard benchmarks without significant overhead in runtime.


   We adopt
  the ground truth Frobenius norm as a reference point, and observe that despite outperforming competing algorithms, FUGAL remains
  substantially below the best achievable alignment under noise. This gap motivates a systematic investigation of FUGAL's algorithmic
  choices to understand where performance is lost and where tuning can recover it. We additionally compare against FAQ, a widely-used and competitive baseline.



\section{Problem Formulation}

The \emph{graph alignment} problem asks: given two
graphs $G_1=(V_1,E_1)$ and $G_2=(V_2,E_2)$ with $|V_1|=|V_2|=n$, find a
bijection $\sigma:V_1\to V_2$ that minimizes the number of edge disagreements (minimizes graph edit distance). The graphs have adjacency
matrices $A,B\in\R^{n\times n}$ respectively.
Formally, the problem can be formulated as:
\begin{equation}
   \min_{P\in\Perm_n} \; \frob{AP - PB}^2,
  \label{eq:qap}
\end{equation}
where $\Perm_n$ denotes the set of permutation matrices.  This is
an instance of the Quadratic Assignment Problem (QAP), which is NP-hard in
general. This paper is an extension of the FUGAL paper and follows the notations, definitions and theories presented in it.


The FUGAL algorithm~\cite{fugal2022} introduces two improvements over
FAQ solver: (1) a Linear Assignment Problem (LAP) supplement that leverages
simple structural graph features (2)~ a regularizing term that guides the solution towards a
quasi-permutation matrix.

< todo : add fugal objective function?>





\section{Experiment Setup}
\label{sec:setup}

\subsection{Datasets}

We evaluate on five real-world graphs drawn from diverse domains.
\Cref{tab:datasets} summarizes their sizes and structural properties.

\begin{table}[H]
  \centering
  \caption{Summary of benchmark graphs.}
  \label{tab:datasets}
  \begin{tabular}{lrrl}
    \toprule
    Dataset & Nodes & Edges & Type \\
    \midrule
    \texttt{highschool}  & 327  & 5818 & proximity \\
    \texttt{multimagna}  & 1004 & 8323 & biological \\
    \texttt{euroroad}    & 1174 & 1417 & infrastructure \\
    \texttt{voles}       & 712  & 2391 & proximity \\
    \texttt{netscience}  & 379  & 914  & collaborations \\
    \bottomrule
  \end{tabular}
\end{table}

\subsection{Noise Model}

To simulate real-world graph corruption, we generate noisy target graphs by
randomly removing edges from the source graph.  Given a noise level
$p \in \{0, 5, 10, 15, 20\}$ (as a percentage), each edge is independently removed with probability corresponding to p, subject to the constraint that no node
becomes isolated.  A random permutation of node labels is simultaneously
applied to the target graph, so the ground-truth correspondence is known, but
must be recovered. The generated source and target graphs are cached for consistent reproducibility of the results throughout the experiments.

\subsection{Evaluation Metrics}

We report two complementary metrics:

\begin{itemize}
  \item \textbf{Accuracy}: fraction of nodes correctly matched to their
    ground-truth counterparts
  \item \textbf{Excess Frobenius norm}:
    $\frob{A\hat{P} - \hat{P}B}^2 - \frob{AP^* - P^*B}^2$,
    the gap between the alignment cost achieved by the algorithm and the cost
    achieved by the ground-truth permutation $P^*$.  This is a strictly
    harder metric than accuracy, since a method can match many nodes correctly
    while still incurring large edge mismatches if the mismatched nodes have
    high degree.
\end{itemize}

All experiments run three independent trials across different noise levels and report means.


\section{Experiment design with key algorithmic choices}


\subsection{Establishing Baseline: FUGAL vs.\ FAQ with Uniform Initialization}
\label{sec:baseline}

\Cref{tab:baseline} compares FUGAL and FAQ when both use uniform initialization.  FUGAL outperforms FAQ on every dataset and noise level in terms of accuracy and Frobenius norm.

\begin{table}[H]
  \centering
  \caption{Accuracy at noise levels 0, 10, 20 (\%) with uniform initialization.
           }
  \label{tab:baseline}
  \setlength{\tabcolsep}{5pt}
  \begin{tabular}{llccccccc}
    \toprule
    & & \multicolumn{2}{c}{noise = 0} & \multicolumn{2}{c}{noise = 10} & \multicolumn{2}{c}{noise = 20} \\
    \cmidrule(lr){3-4}\cmidrule(lr){5-6}\cmidrule(lr){7-8}
    Dataset & & FUGAL & FAQ & FUGAL & FAQ & FUGAL & FAQ \\
    \midrule
    highschool  & acc  & \textbf{1.000} & \textbf{1.000} & \textbf{1.000} & 0.692 & \textbf{1.000} & 0.284 \\
    multimanga  & acc  & \textbf{0.846} & 0.482 & \textbf{0.789} & 0.352 & \textbf{0.635} & 0.291 \\
    euroroad    & acc  & \textbf{0.944} & 0.013 & \textbf{0.732} & 0.005 & \textbf{0.248} & 0.005 \\
    voles       & acc  & \textbf{0.948} & 0.691 & \textbf{0.936} & 0.163 & \textbf{0.893} & 0.039 \\
    netscience  & acc  & \textbf{0.685} & 0.317 & \textbf{0.641} & 0.084 & \textbf{0.546} & 0.044 \\
    \bottomrule
  \end{tabular}
\end{table}



\subsection{Effect of warm-start initialization}
\label{sec:init}


We next apply warm-start initialization,  a computationally expensive but high-quality starting point to both algorithms. We set the baselines with FUGAL and FAQ, with and without the warm-start initializations. FAQ without initialization performed poorly, and hence, we are only looking at fugal\_init, faq\_init and fugal for our comparisons.

faq\_init gives the best quality alignment for netscience, euroroad and voles datasets.

\begin{table}[H]
  \centering
  \caption{Accuracy at noise = 10 and 20 — comparing uniform vs.\ warm-start initialization.
           \texttt{Fugal\_init} best accuracy across $\mu$ values.}
  \label{tab:init}
  \setlength{\tabcolsep}{5pt}
  \begin{tabular}{lcccccccc}
    \toprule
    & \multicolumn{4}{c}{noise = 10} & \multicolumn{4}{c}{noise = 20} \\
    \cmidrule(lr){2-5}\cmidrule(lr){6-9}
    Dataset & FUGAL & FUGAL$_{\text{init}}$ & FAQ & FAQ$_{\text{init}}$ &
              FUGAL & FUGAL$_{\text{init}}$ & FAQ & FAQ$_{\text{init}}$ \\
    \midrule
    highschool & 1.000 & 1.000 & 0.692 & 1.000 & 1.000 & 1.000 & 0.284 & 1.000 \\
    multimanga & 0.789 & 0.759 & 0.352 & 0.715 & 0.635 & 0.620 & 0.291 & 0.568 \\
    euroroad   & 0.732 & 0.746 & 0.005 & 0.819 & 0.248 & 0.274 & 0.005 & 0.594 \\
    voles      & 0.936 & 0.943 & 0.163 & 0.942 & 0.893 & 0.905 & 0.039 & 0.908 \\
    netscience & 0.641 & 0.672 & 0.084 & 0.669 & 0.546 & 0.573 & 0.044 & 0.550 \\
    \bottomrule
  \end{tabular}
\end{table}

Two observations stand out:

\begin{enumerate}
  \item \textbf{FUGAL gains little from warm-start.}  Initialization improves FUGAL's accuracy by at most 2--4 percentage points across datasets.  FUGAL's Frank-Wolfe iterations appear to escape poor initializations effectively on their own.

  \item \textbf{FAQ gains enormously from warm-start.}  FAQ's accuracy jumps dramatically. It also outperforms FUGAL in 3 of the 5 datasets. However, this improvement comes at significant computational cost: warm-start initialization requires expensive seed-alignment computation.
\end{enumerate}

Despite the small benefit, FUGAL with warm-start initialization provides a useful starting point for hyperparameter fine-tuning, studied next.

\subsection{Hyperparameter grid search for fugal\_init}

FUGAL has two key hyperparameters: the feature weight $\muval$ and the regularization step $\lams$.  We perform a grid search over $\muval \in \{0.05, 0.1, 0.3, 0.5, 1.0, 2.0, 3.0\}$ and $\lams \in \{0.2, 0.4, 0.6, 0.8, 1.0\}$ for each dataset and noise level.
\\



\begin{figure}[H]
    \centering
    \includegraphics[width=0.8\linewidth]{gridsearch_netscience.png}
    \caption{Hyperparameter grid search on netscience dataset for fugal\_init. Excess frobenius norm is captured for each combination. fugal\_init is getting equal or better results than faq\_init. Circled section shows instances where fugal\_init outperforms faq\_init. }
    \label{fig:placeholder}
\end{figure}



Further, gradually reducing $\lams$ toward the
  end of the optimization yielded consistent improvements.


\Cref{tab:finetune} compares the best fine-tuned FUGAL$_{\text{init}}$ against FAQ$_{\text{init}}$.

\begin{table}[H]
  \centering
  \caption{Accuracy — fine-tuned FUGAL$_{\text{init}}$ vs.\ FAQ$_{\text{init}}$.
           Bold indicates the better result. $\uparrow$ FUGAL wins, $\downarrow$ FAQ wins.}
  \label{tab:finetune}
  \setlength{\tabcolsep}{4pt}
  \begin{tabular}{lcccccc}
    \toprule
    & \multicolumn{2}{c}{noise = 5} & \multicolumn{2}{c}{noise = 10 (if avail)} & \multicolumn{2}{c}{noise = 20} \\
    \cmidrule(lr){2-3}\cmidrule(lr){4-5}\cmidrule(lr){6-7}
    Dataset & FUGAL$_{\text{init}}^*$ & FAQ$_{\text{init}}$ & FUGAL$_{\text{init}}^*$ & FAQ$_{\text{init}}$ &
              FUGAL$_{\text{init}}^*$ & FAQ$_{\text{init}}$ \\
    \midrule
    highschool & \textbf{1.000} & \textbf{1.000} & \textbf{1.000} & \textbf{1.000} & \textbf{1.000} & \textbf{1.000} \\
    multimanga & \textbf{0.833} $\uparrow$ & 0.821 & \textbf{0.758} $\uparrow$ & 0.715 & \textbf{0.629} $\uparrow$ & 0.567 \\
    voles      & \textbf{0.950} $\uparrow$ & 0.946 & \textbf{0.943} $\uparrow$ & 0.942 & \textbf{0.916} $\uparrow$ & 0.908 \\
    netscience & 0.690 & \textbf{0.691} & \textbf{0.683} $\uparrow$ & 0.669 & \textbf{0.585} $\uparrow$ & 0.550 \\
    euroroad   & 0.868 & \textbf{0.894} $\downarrow$ & --- & \textbf{0.819} $\downarrow$ & 0.356 & \textbf{0.594} $\downarrow$ \\
    \bottomrule
  \end{tabular}
\end{table}

This experiment gives us the best mu, lam\_step values per dataset noise combinations, which we use for the future experiments.
Fine-tuned FUGAL$_{\text{init}}$ showed improvement in getting closer to the ground truth alignment in 4 of the 5 datasets. The lone exception is \texttt{euroroad}, where FAQ$_{\text{init}}$ holds a substantial lead.  We investigate this case separately in \Cref{sec:euroroad}.

\subsection{The EuroRoad Case: Learning Rate Schedule}
\label{sec:euroroad}

\texttt{EuroRoad} is the sparsest graph in our benchmark. The wide gap between fine-tuned FUGAL$_{\text{init}}$ and FAQ$_{\text{init}}$ prompted further investigation and identified that FUGAL could benefit from the exact line search learning schedule, like in FAQ. Currently,
FUGAL uses a fixed step size for its learning rate.

We adopted FAQ's exact line search schedule within FUGAL$_{\text{init}}$ \Cref{tab:newalpha} shows the improvement:

\begin{table}[H]
  \centering
  \caption{Accuracy on EuroRoad — effect of new learning rate schedule.}
  \label{tab:newalpha}
  \begin{tabular}{lcccc}
    \toprule
    Noise & FUGAL$_{\text{init}}$ (fixed step) & FUGAL$_{\text{init}}$ (newalpha) & FAQ$_{\text{init}}$ \\
    \midrule
    5  & 0.868 & \textbf{0.890} & 0.894 \\
    10 & --- & \textbf{0.821} & 0.819 \\
    15 & 0.558 & \textbf{0.677} & 0.741 \\
    20 & 0.356 & \textbf{0.476} & 0.594 \\
    \bottomrule
  \end{tabular}
\end{table}

The new learning schedule substantially narrows the gap with FAQ$_{\text{init}}$ at all noise levels, and even surpasses it at noise 10.

\subsection{Projection Matrix}
FUGAL uses sinkhorn while FAQ uses hungarian. Soft projection via Sinkhorn is yielding optimal results in 4 of the 5 datasets, while in euroroad, hungarian seems to outperform sinkhorn. This is an interesting observation as it shows that sinkhorn is not always the best choice out there.


\section{Results}



\begin{figure}[H]
    \centering
    \includegraphics[width=0.5\linewidth]{frob_gap_euroroad.png}
\end{figure}

\begin{figure}[H]
    \centering
    \includegraphics[width=0.5\linewidth]{frob_gap_highschool.png}
\end{figure}

\begin{figure}[H]
    \centering
    \includegraphics[width=0.5\linewidth]{frob_gap_netscience.png}
\end{figure}

\begin{figure}[H]
    \centering
    \includegraphics[width=0.5\linewidth]{frob_gap_multimanga.png}
\end{figure}

\begin{figure}[H]
    \centering
    \includegraphics[width=0.5\linewidth]{frob_gap_voles.png}
\end{figure}
\section{Discussion}
\label{sec:discussion}

\paragraph{Why FUGAL gains little from initialization.}
FUGAL's Frank-Wolfe iterations, augmented by the LAP term, appear to guide the solution towards good alignment even from a uniform start.

\paragraph{Why FAQ gains enormously from initialization.}
FAQ relies solely on the quadratic objective and has no structural bias.  From a random start, it tends to converge to poor local optima.  With a good initialization, FAQ converges to high-quality solutions.  This is consistent with the broader literature on non-convex optimization for graph matching.



\paragraph{Frobenius norm vs.\ accuracy.}
Excess Frobenius is a strictly harder metric than accuracy.  A method can achieve high node-matching accuracy (e.g., 83\% of nodes correctly identified) while still incurring large edge mismatches if the mismatched nodes have high degrees. We observe this most clearly on multimanga and netscience, where algorithms may differ on accuracy by a small margin but differ substantially on excess Frobenius.

\paragraph{$\muval$ sensitivity.}
The optimal $\muval$ is dataset- and noise-dependent.  On dense graphs (highschool, multimanga), small values ($\muval \approx 0.05$--$0.5$) suffice, while on sparser graphs (euroroad, netscience at higher noise), larger values ($\muval \approx 1.0$--$2.0$) are preferred.

\paragraph{$\lams$ sensitivity.}
Optimal $\lams$ values cluster around $0.2$--$0.4$ for most datasets, indicating that a finer regularization step generally performs better.

\paragraph{Future experiments.}
A natural next step is to evaluate FUGAL \emph{without} warm-start but \emph{with} the exact line search schedule.  If this combination closes the gap with warm-start results, it would offer a compelling alignment-runtime tradeoff: competitive accuracy at the cost of only the FUGAL runtime (no expensive seed computation).  Adaptive or learned hyperparameter selection (e.g., using graph structural statistics to predict good $\muval$ and $\lams$) is another promising direction.

\section{Conclusion}
\label{sec:conclusion}

We systematically studied four algorithmic dimensions for graph alignment problem: initialization, learning rate schedule, projection type, and objective terms, across five real-world graphs at five noise levels.
From the different experiments we conducted, we find that finding a single universal algorithm that works optimally on all the datasets is hard, and there is no clear winner. But, on a positive note, we discover a few promising algorithmic trends.

Our main findings are:

\begin{enumerate}
  \item \textbf{FUGAL outperforms FAQ with uniform initialization} on all five datasets, confirming prior results.

  \item \textbf{FUGAL is robust to initialization}; FAQ is not.  Warm-start initialization yields dramatic improvements for FAQ but only marginal gains for FUGAL.

  \item \textbf{Fine-tuned FUGAL$_{\text{init}}$ beats FAQ$_{\text{init}}$} on four of five datasets at moderate-to-high noise, with the exception of the extremely sparse EuroRoad network.

  \item \textbf{The learning rate schedule is the dominant factor for sparse graphs.}  Adopting FAQ's exact line search within FUGAL narrows the EuroRoad gap substantially and even surpasses FAQ$_{\text{init}}$ at noise level 10.

  \item \textbf{Optimal hyperparameters are graph-dependent.}  Dense, feature-rich graphs prefer small $\muval$ and moderate $\lams$; sparse graphs under heavy noise benefit from larger $\muval$.
\end{enumerate}

Warm-start initialization is computationally expensive and unlikely to scale to large graphs.  Whether FUGAL with uniform initialization and the new learning rate schedule can approach warm-start quality remains an open question and the most promising direction for future work.



\begin{thebibliography}{9}


\bibitem{vogelstein2015}
J.~T.~Vogelstein et al.
\newblock Fast approximate quadratic programming for graph matching.
\newblock \emph{PLOS ONE}, 10(4):e0121002, 2015.

\end{thebibliography}

\end{document}